% from ega4-9-fr.tex
\oldpage[IV-3]{74}

\begin{enumerate}[label=\upshape{II)}]
  \item Soit $\mathcal L$ un $\mathcal O_X$-Module inversible. Alors l'ensemble
  $E'$ des $s\in S$ tels que $\mathcal L_s$ soit un $\mathcal O_{X_s}$-Module
  ample (resp. très ample) relativement à $f_s$, est ind-constructible dans
  $S$.
\end{enumerate}
\end{proposition}

I) Vérifions les conditions (i) et (ii) de (9.2.1). En ce qui concerne la
condition (9.2.1, (i)), elle résulte, pour les propriétés (i) et (ii), de
(2.7.1, (xiii) et (xiv))~; pour les propriétés (iii) et (iv), elle résulte de
(9.1.5). Vérifions ensuite la condition (9.2.1, (ii)), en supposant donc $S$
noethérien, intègre et de point générique $\eta$, et que
$f_\eta:X_\eta\to Y_\eta$ ait l'une des propriétés (i) à (iv) de l'énoncé.
Appliquant (8.10.5, (viii), (ix), (xiii) et (xiv)) suivant la méthode de
(8.1.2, \emph{a)}), on voit d'abord qu'il existe un voisinage ouvert $U$ de
$\eta$ tel que, si $V$ et $W$ sont les images réciproques de $U$ dans $X$ et
$Y$ respectivement par les morphismes structuraux, la restriction $V\to W$ de
$f$ a celle des propriétés (i) à (iv) que l'on considère. La conclusion
résulte alors de ce que ces propriétés sont toutes stables par changement de
base.

II) On procède de la même façon. La condition (9.2.1, (i)) résulte cette fois
de (2.7.2). Pour la condition (9.2.1, (ii)), avec les mêmes notations que dans
I), il résulte de (8.10.5.2) que pour un voisinage $U$ de $\eta$ dans $S$, la
restriction $\mathcal L|V$ est ample (resp. très ample) relativement à la
restriction $V\to W$ de $f$. La conclusion résulte encore de la stabilité des
propriétés considérées par changement de base (\textbf{II}, 4.6.13 et
4.4.10).

On peut améliorer (9.6.2, II) sous certaines conditions~:

\begin{proposition}[9.6.3]
\label{IV.9.6.3-fr}
Soient $X$, $Y$ deux $S$-préschémas de présentation finie, $f:X\to Y$ un
$S$-morphisme propre, $\mathcal L$ un $\mathcal O_X$-Module inversible. Alors
l'ensemble $E$ (resp. $E'$) des $s\in S$ tels que $\mathcal L_s$ soit un
$\mathcal O_{X_s}$-Module ample (resp. très ample) relativement à
$f_s:X_s\to Y_s$ est localement constructible dans $S$.
\end{proposition}

Soient $k$ un corps, $Z$, $T$ deux préschémas algébriques sur $k$, $\mathcal M$
un $\mathcal O_Z$-Module inversible, $g:Z\to T$ un $k$-morphisme de type fini~;
alors, si $\mathbf P(Z,T,\mathcal M,g,k)$ désigne la relation~: «~$\mathcal M$ est ample
(resp. très ample) relativement à $g$~», on a déjà remarqué dans (9.6.2) que
$\mathbf P(Z,T,\mathcal M,g,k)$ vérifie la condition (9.2.1, (i)) en vertu de (2.7.2).
On sait déjà d'autre part que $E$ et $E'$ sont ind-constructibles. Il reste donc
à voir que si $S$ est noethérien, intègre, de point générique $\eta$ et si
$\eta\in S-E$ (resp. $\eta\in S-E'$), alors $S-E$ (resp. $S-E'$) contient un
voisinage de $\eta$. Nous considérerons séparément le cas de $E'$ et celui de
$E$.

I) \emph{Cas de $E'$.} --- Notons que puisque $f$ est séparé et $Y$
quasi-compact, il existe un entier $h$ tel que, pour $q>h$, on ait
$R^qf_*(\mathcal L)=0$ (\textbf{III}, 1.4.12)~; d'autre part, puisque $f$ est
propre et $Y$ noethérien, les $R^qf_*(\mathcal L)$ sont tous des
$\mathcal O_Y$-Modules cohérents (\textbf{III}, 3.2.1)~; comme ils sont nuls
sauf pour un nombre fini de valeurs de $q$, le théorème de platitude générique
(6.9.1) montre qu'en restreignant $S$ à un voisinage de $\eta$, on peut
supposer que $\mathcal L$ et les $R^qf_*(\mathcal L)$ sont tous $S$-plats. On
conclut alors de (\textbf{III}, 6.9.9) que l'homomorphisme canonique
\begin{equation}
\label{IV.9.6.3.1-fr}
 f_*(\mathcal L)\otimes_{\mathcal O_S} k(s)
 \longrightarrow (f_s)_*(\mathcal L_s)
 \tag{9.6.3.1}
\end{equation}
est un \emph{isomorphisme}.

